\chapter{Related Work}\label{chap:relatedWork}

\section{NIMEs}
TODO: A quick explanation of NIMEs and their history in research. Talk specifically about the distinction of real-time interactive NIMEs as apposed to other NIMEs, and 

\subsection{Symbolic Approaches}

\subsection{Machine Learning Approaches}

\section{RNN Prediction Trees}
Similar approaches of using RNNs to generate nodes in prediction trees, which can then be searched using heuristic based methods, such as MCTS, are not present in the state-of-the-art for computer music. However, similar approaches are state-of-the-art techniques in other computing fields, particular in path planning for dynamic environments. For this problem, regular search techniques such as A* cannot generate effective paths, as during the execution of the path the environment around the agent will be dynamically changing in a complex multi-dimensional way, which can unpredictably affect the validity and optimality of generated paths. Hence, because the states after executing an action are non-deterministic, RNNs are used to generate predicted nodes for each point in the search tree. These RNNs are good at accounting for the complex dynamic environment, and can provide reasonable predictions for what some realistic future scenarios could be after executing a set of actions. The agent can then search through these RNN created environmental prediction trees, often using MCTS, to find paths which more effectively account for the unpredictable environment. (cite specific cases of these articles, and their results, here). The use of very similar techniques provides an analogy to our IMPS approach. Similarly to a dynamic physical environment that an agent is trying to pathfind within, in music we often have broad structural goals we want to adhere to, or important musical moments, such as a chorus, that we set a goal to build toward. However, in the short term we have to adapt to the complex multi-dimensional and hard to symbolically encode rules of what makes music sound good. Similarly to path planning in dynamic environments, we are aiming to use RNNs to predict and navigate the flow of music on a small scale, while search algorithms and symbolic heuristics control the broader cohesiveness and longer-term structural goals of a performance. We also see that current IMPS methods face analagous issues to only using RNNs or only using symbolic tree search in path finding in dynamic environments. When only symbolic reasoning is used, the agent will maintain some form of of goal-orientation and structure, but on a smaller scale it will take very inefficient paths as it cannot accurately predict its environment and must respond only when things go so wrong that its symbolic heuristics pick up on the issue. Similarly only using RNNs would lead to an agent which can very smoothly and effectively navigate the environment on a small scale, but would meander aimlessly without a strong long term direction. Hence, the techniques used in these papers are likely to be helpful in guiding the approaches we take in our paper, and so we will now analyse them further.

TODO: Look into depth on the papers in the field, what things worked, what didn't, what was the effect of different methodologies and approaches.
